\documentclass[aps,rmp,reprint,amsmath,amssymb,graphicx,longbibliography]{revtex4-1}

\usepackage[utf8]{inputenc}
\usepackage{bm}
\usepackage{graphicx}
\usepackage{booktabs}
\usepackage{float}
\usepackage{listings}
\usepackage{xcolor}
\usepackage{hyperref}
\usepackage{placeins}

\hypersetup{
    breaklinks=true,
    colorlinks=true,
    linkcolor=blue,
    citecolor=blue,
    urlcolor=blue
}
\setcitestyle{numbers,square}

\begin{document}

\title{Project 2: Neural Networks for Regression and Classification: Implementation, Optimization, and Critical Evaluation}

\author{Teresa Macarulla Uriarte}
\affiliation{FYS-STK3155 Applied Data Analytics and Machine Learning}
\affiliation{University of Oslo}
\date{\today}

\begin{abstract}
We develop a feed-forward neural-network implementation in NumPy for regression and
multiclass classification, with manually coded backpropagation and mini-batch
optimization. Regression on noisy Runge-function data is compared with polynomial
ordinary least squares (OLS), Ridge and Lasso; the influence of optimizers,
activation functions, network size and regularization is examined. The lowest
observed held-out regression error is obtained with an unregularized ReLU network
(test MSE $0.008492$), compared with $0.010240$ for degree-14 OLS, although the
polynomial model is substantially simpler. For classification, a softmax output
and cross-entropy loss are applied to the full MNIST handwritten-digit dataset.
The validation-selected $L_2$-regularized network reaches $97.90\%$ accuracy
and cross-entropy $0.064215$ on the official test set; an unregularized network
reaches $97.87\%$, versus $92.45\%$ for multinomial logistic regression.
The results illustrate both the advantages of learned nonlinear features and
the importance of validation, numerical gradient checks and appropriately
matched model complexity.
\end{abstract}

\maketitle

\tableofcontents

% ============================================================
\section{Introduction}

Regression and classification are central supervised-learning tasks, but their
outputs and optimization objectives differ. Regression estimates continuous
responses, whereas classification assigns observations to discrete categories.
Neural networks offer a common framework for both: successive trainable layers
learn nonlinear representations, while the output activation and loss function
are chosen to match the task. Their flexibility must nevertheless be weighed
against computational cost, sensitivity to optimization settings and the risk
of overfitting \cite{Book_Ch8}.

The purpose of this project is to implement a feed-forward neural network (NN)
and backpropagation from scratch, building on the gradient-based optimization
methods studied in Project~1 \cite{Project_1}. First, noisy observations of
Runge's function are used to compare neural-network regression with polynomial
ordinary least squares (OLS), Ridge and Lasso. We examine the effects of
optimizer, learning rate, hidden-layer activation, width, depth and weight
regularization. Next, the same code is adapted to multiclass classification
with a softmax output and cross-entropy loss. On the full MNIST handwritten-
digit dataset \cite{lecun2010mnist}, we compare a multinomial logistic-
regression baseline with networks containing hidden layers, explore their
hyperparameters, and evaluate a validation-selected classifier on a held-out
test set.

The implementation uses three Python files: \texttt{main.py} controls the
experiments, \texttt{utilities.py} contains the numerical methods and training
routines, and \texttt{plot\_generator.py} generates the figures. The source
code is available in the
\href{https://github.uio.no/teresm/Machine-Learning/tree/main/P2}{project repository}.
The following sections present the mathematical formulation, implementation,
results and a critical assessment of the methods for the two datasets.


% ============================================================
\section{Theory}


\subsection{Mean-squared error and regularization}
\label{subsec:mse_gradients}

For regression, the mean-squared error (MSE) quantifies the average squared
residual between the target values $\mathbf{y}$ and predictions
$\hat{\mathbf{y}}$. For a linear model with design matrix
$X\in\mathbb{R}^{n\times p}$ and parameter vector
$\boldsymbol{\theta}\in\mathbb{R}^{p}$, the predictions are
$\hat{\mathbf{y}}=X\boldsymbol{\theta}$ and the cost is \cite{Book_Ch3}
\begin{equation}
    C_{\mathrm{MSE}}(\boldsymbol{\theta})
    =\frac{1}{n}\sum_{i=1}^{n}(y_i-\hat y_i)^2
    =\frac{1}{n}\|\mathbf{y}-X\boldsymbol{\theta}\|_2^2.
    \label{eq:MSE}
\end{equation}
Differentiating with respect to the parameters gives
\begin{equation}
    \nabla_{\boldsymbol{\theta}}C_{\mathrm{MSE}}
    =\frac{2}{n}X^{\mathsf T}(X\boldsymbol{\theta}-\mathbf{y}).
\end{equation}

To control model complexity, the cost can be augmented with either an
$L_1$ (Lasso) or squared $L_2$ (Ridge) penalty:
\begin{equation}
    C_{L_1}=C_{\mathrm{MSE}}+\lambda\sum_{j\in\mathcal{P}}|\theta_j|,
    \qquad
    C_{L_2}=C_{\mathrm{MSE}}+\lambda\sum_{j\in\mathcal{P}}\theta_j^2,
\end{equation}
where $\lambda\geq0$ controls the regularization strength and
$\mathcal{P}$ denotes the penalized parameters (usually excluding the
intercept) \cite{Book_Ch3}. Their gradients are
\begin{align}
    \partial_{\boldsymbol{\theta}}C_{L_1}
    &\ni \frac{2}{n}X^{\mathsf T}(X\boldsymbol{\theta}-\mathbf{y})
      +\lambda\mathbf{s},\\
    \nabla_{\boldsymbol{\theta}}C_{L_2}
    &=\frac{2}{n}X^{\mathsf T}(X\boldsymbol{\theta}-\mathbf{y})
      +2\lambda\boldsymbol{\theta}_{\mathcal{P}},
\end{align}
where $s_j=\operatorname{sign}(\theta_j)$ for penalized nonzero
coefficients, $s_j\in[-1,1]$ at zero, and $s_j=0$ for excluded
coefficients. The vector $\boldsymbol{\theta}_{\mathcal{P}}$ equals
$\boldsymbol{\theta}$ on penalized entries and zero elsewhere.
Thus, the $L_1$ expression is a subgradient at zero, whereas the squared
$L_2$ penalty is differentiable everywhere. We use the convention
$1/n$ (rather than $1/(2n)$) to match the MSE implementation.


\subsection{Binary cross-entropy and regularization}
\label{subsec:binary_cross_entropy}

For binary classification, each target $y_i\in\{0,1\}$ is modeled by a
Bernoulli probability $p_i=\sigma(z_i)$, where
$z_i=\mathbf{x}_i^{\mathsf T}\boldsymbol{\theta}$ and
$\sigma(z)=1/(1+e^{-z})$ \cite{Book_Ch5}. The average binary cross-entropy (log loss) is
\begin{equation}
 C_{\mathrm{BCE}}(\boldsymbol{\theta})=
 -\frac{1}{n}\sum_{i=1}^{n}
 \left[y_i\ln p_i+(1-y_i)\ln(1-p_i)\right].
\end{equation}
Since $\partial C_{\mathrm{BCE}}/\partial z_i=(p_i-y_i)/n$,
the chain rule gives
\begin{equation}
 \nabla_{\boldsymbol{\theta}}C_{\mathrm{BCE}}
 =\frac{1}{n}X^{\mathsf T}(\mathbf p-\mathbf y).
\end{equation}
Adding $L_1$ or squared $L_2$ regularization over the same penalized
parameter set $\mathcal{P}$ as in the MSE case gives
\begin{equation}
 C_{\mathrm{BCE},L_1}=C_{\mathrm{BCE}}
 +\lambda\sum_{j\in\mathcal{P}}|\theta_j|,\qquad
 C_{\mathrm{BCE},L_2}=C_{\mathrm{BCE}}
 +\lambda\sum_{j\in\mathcal{P}}\theta_j^2.
\end{equation}
Their corresponding (sub)gradients are
\begin{align}
 \partial_{\boldsymbol{\theta}}C_{\mathrm{BCE},L_1}
 &\ni \frac{1}{n}X^{\mathsf T}(\mathbf p-\mathbf y)
 +\lambda\mathbf s,\\
 \nabla_{\boldsymbol{\theta}}C_{\mathrm{BCE},L_2}
 &= \frac{1}{n}X^{\mathsf T}(\mathbf p-\mathbf y)
 +2\lambda\boldsymbol{\theta}_{\mathcal{P}}.
\end{align}
Here $\mathbf s$ and $\boldsymbol{\theta}_{\mathcal{P}}$ follow the
conventions introduced for MSE regularization \cite{Book_Ch5}.


\subsection{Multiclass cross-entropy and softmax}
\label{subsec:softmax_cross_entropy}

Multiclass cross-entropy extends binary classification to $K$ classes.
Instead of a binary target $y_i\in\{0,1\}$, we use one-hot encoding,
where $y_{ik}=1$ if observation $i$ belongs to class $k$, and zero
otherwise. The softmax function provides the predicted probability
$p(C=k\mid\mathbf{x}_i)$ for each class,

\begin{equation}
    p(C=k\mid\mathbf{x})=\frac{\mathbf{exp}(\mathbf{x}^T\,\boldsymbol{\theta}_k)}{\sum^{K-1}_{r=0}\mathbf{exp}(\mathbf{x}^T\,\boldsymbol{\theta}_r)}.
\end{equation}

The multiclass cross-entropy cost function is
\begin{equation}
    C_{\mathrm{CE}} =
    -\frac{1}{n}\sum_{i=1}^{n}\sum_{k=0}^{K-1}
    y_{ik}\ln p(C=k\mid\mathbf{x}_i).
\end{equation}

Its gradient with respect to the parameters of class $k$ is
\begin{equation}
    \frac{\partial C_{\mathrm{CE}}}
    {\partial\boldsymbol{\theta}_k}
    =-\frac{1}{n}X^{\mathsf T}
    (\mathbf{y}_k-\mathbf{p}_k),
\end{equation}
where $\boldsymbol{\theta}_k$ denotes the parameter vector for class
$k$, while $\mathbf{y}_k$ and $\mathbf{p}_k$ contain the corresponding
target indicators and predicted probabilities.

For Neural Networks (NNs), the same loss functions are evaluated using
predictions obtained through successive network layers. The cost
function is minimized with respect to all network weights and biases,
and its gradients are calculated using backpropagation, as discussed
in Chapter 8 of the lecture notes \cite{Book_Ch8}. In this project, MSE is used for
regression, while multiclass cross-entropy is used for classification.


\subsection{Activation functions and their derivatives}
\label{subsec:activation_functions}

Activation functions introduce nonlinearity into NNs, allowing
them to learn relationships beyond linear combinations of the inputs.
Their first derivatives are required during backpropagation to compute
the gradients of the cost function.

The \textit{sigmoid} activation function maps its input to the interval
$(0,1)$ \cite{Book_Ch8} and is defined as
\begin{equation}
    \sigma(z)=\frac{1}{1+e^{-z}},
    \qquad
    \sigma'(z)=\sigma(z)[1-\sigma(z)].
\end{equation}
The sigmoid saturates for large positive or negative inputs, potentially
leading to vanishing gradients.

The \textit{Rectified Linear Unit} (ReLU) is defined as
\begin{equation}
    \mathrm{ReLU}(z)=\max(0,z),
    \qquad
    \mathrm{ReLU}'(z)=
    \begin{cases}
        1, & z>0,\\
        0, & z<0.
    \end{cases}
\end{equation}
Unlike the sigmoid, ReLU does not saturate for positive inputs \cite{Book_Ch8}.
However, its zero gradient for negative inputs can prevent certain
neurons from updating their weights.

The \textit{Leaky ReLU} introduces a small nonzero slope for negative
inputs, controlled by a parameter $\alpha$, typically $\alpha=0.01$:
\begin{equation}
    \mathrm{LReLU}(z)=
    \begin{cases}
        z, & z>0,\\
        \alpha z, & z\leq 0,
    \end{cases}
    \qquad
    \mathrm{LReLU}'(z)=
    \begin{cases}
        1, & z>0,\\
        \alpha, & z<0.
    \end{cases}
\end{equation}
This modification allows gradients to propagate even when the
pre-activation values are negative \cite{Book_Ch8}. Both ReLU and Leaky ReLU are
not differentiable at $z=0$, where a convention must be chosen
for numerical implementation.


\subsection{Feed-forward regression and backpropagation}
\label{subsec:ffnn_backprop}

A feed-forward NN represents a nonlinear mapping by composing
successive affine transformations and activation functions. For a batch of
$m$ observations stored in rows, the forward pass through layer $\ell$ is
\begin{align}
 Z^{(\ell)} &= A^{(\ell-1)}W^{(\ell)}
                 +\mathbf{1}(\mathbf{b}^{(\ell)})^{\mathsf T},\\
 A^{(\ell)} &= g_{\ell}\!\left(Z^{(\ell)}\right),
 \qquad A^{(0)}=X,
\end{align}
where $W^{(\ell)}$ and $\mathbf{b}^{(\ell)}$ are the layer weights and
biases. In the regression networks considered here, hidden layers use
sigmoid, ReLU or Leaky ReLU activation, while the final layer uses the identity function,
$g_L(z)=z$, so that the predicted response is not restricted to $(0,1)$.

The gradients of the unregularized MSE are obtained by backpropagating
errors through the network using the chain rule \cite{Book_Ch8}. 
For a batch of $m$ observations, the unregularized MSE is
\begin{equation}
    C_{\mathcal{B}} =
    \frac{1}{m}\sum_{i \in \mathcal{B}}
    (y_i-\hat{y}_i)^2.
\end{equation}
Here, $m$ denotes the batch size, while $n$ in Eq.~\eqref{eq:MSE}
denotes the total number of observations. For a linear output activation,
the output-layer error is
\begin{equation}
    \delta^{(L)}
    = \frac{2}{m}(\hat{\mathbf{y}}-\mathbf{y}).
\end{equation}

For a hidden layer, the errors satisfy
\begin{equation}
 \delta^{(\ell)}=\left[\delta^{(\ell+1)}
   (W^{(\ell+1)})^{\mathsf T}\right]
   \odot g_\ell'\!\left(Z^{(\ell)}\right),
\end{equation}
where $\odot$ denotes element-wise multiplication \cite{Book_Ch8}. The corresponding
parameter gradients are
\begin{equation}
 \frac{\partial C}{\partial W^{(\ell)}}
   =(A^{(\ell-1)})^{\mathsf T}\delta^{(\ell)},
   \label{eq:parameter_gradient}
\end{equation}
\begin{equation}
 \frac{\partial C}{\partial \mathbf b^{(\ell)}}
   =(\delta^{(\ell)})^{\mathsf T}\mathbf{1}.
\end{equation}
The output error contains the factor of two because the cost uses MSE,
not half the MSE. These gradients can be combined with full-batch or
mini-batch parameter updates.


\subsection{Regularization of neural-network weights}
\label{subsec:nn_regularization_theory}

In the feed-forward networks, the same $L_1$ and squared $L_2$
penalties introduced above are applied to the weight matrices of
all layers, excluding biases. For a mini-batch $\mathcal{B}$ of
$m$ observations, the regularized objective is
\begin{align}
 C_{L_1} &= C_{\mathcal{B}}+
   \lambda\sum_{\ell=1}^{L}\sum_{j,k}|W_{jk}^{(\ell)}|,\\
 C_{L_2} &= C_{\mathcal{B}}+
   \lambda\sum_{\ell=1}^{L}\sum_{j,k}(W_{jk}^{(\ell)})^2,
\end{align}
where $C_{\mathcal{B}}$ is the unregularized batch MSE.
Accordingly, the weight gradients (Eq.~\eqref{eq:parameter_gradient}) from backpropagation acquire an
additional term $\lambda\,\mathrm{sign}(W^{(\ell)})$ for $L_1$, or
$2\lambda W^{(\ell)}$ for squared $L_2$; the bias gradients remain
unchanged. For $L_1$ a subgradient of zero is chosen at a zero weight.
The penalized objective is optimized during training, whereas plain MSE
is retained as the common measure of predictive performance
\cite{Book_Ch8}.


\subsection{Multiclass network output and evaluation metrics}
\label{subsec:mnist_theory}

The softmax and cross-entropy formulation of
Sec.~\ref{subsec:softmax_cross_entropy} also applies to a feed-forward
classifier: the final-layer pre-activations $Z^{(L)}$ replace the linear
scores of multinomial logistic regression. For a mini-batch of $m$
observations, let $P$ contain the predicted class probabilities and
$Y$ their one-hot targets. The combination of softmax and cross-entropy
simplifies the output-layer error to \cite{Book_Ch5,Book_Ch8}
\begin{equation}
    \delta^{(L)}_{\mathrm{CE}}
    =\frac{\partial C_{\mathrm{CE}}}{\partial Z^{(L)}}
    =\frac{1}{m}(P-Y).
    \label{eq:mnist_output_delta}
\end{equation}
The hidden-layer recursion and parameter gradients of
Sec.~\ref{subsec:ffnn_backprop} remain unchanged; no separate softmax
derivative is needed in backpropagation. For $L_1$ or squared $L_2$
classification, the MSE term in Sec.~\ref{subsec:nn_regularization_theory}
is replaced by cross-entropy, while the weight-penalty terms are unchanged.
A network with no hidden layer therefore reduces to multinomial logistic
regression.

For an observation with true digit $t_i$, the predicted label is
$\hat t_i=\arg\max_k p_{ik}$. Classification accuracy is
\begin{equation}
    \mathrm{Accuracy}=
    \frac{1}{n}\sum_{i=1}^{n}\mathbb{I}(t_i=\hat t_i),
    \label{eq:mnist_accuracy}
\end{equation}
where $\mathbb{I}$ is the indicator function. Unlike accuracy, cross-entropy
also reflects the probability assigned to the correct class. To examine
errors by digit, the confusion matrix is defined as
\begin{equation}
    M_{jk}=\sum_{i=1}^{n}
    \mathbb{I}(t_i=j)\,\mathbb{I}(\hat t_i=k),
    \label{eq:mnist_confusion}
\end{equation}
with true labels along rows and predicted labels along columns. For
class $k$, precision is $M_{kk}/\sum_j M_{jk}$, recall is
$M_{kk}/\sum_j M_{kj}$, and $F_1$ is their harmonic mean
\cite{Book_Ch5}.


% ============================================================
\section{Methods and Implementation}

\subsection{Runge data, preprocessing, and network architecture}
\label{subsec:partb_methods}

The regression study used $n=200$ noisy observations of Runge's function,
$f(x)=(1+25x^2)^{-1}$, with inputs sampled uniformly on $[-1,1]$ and
additive Gaussian noise of standard deviation $\sigma=0.1$. The data
were generated with NumPy \cite{harris2020array} seed 2026, reproducing the dataset used in
Project~1 \cite{Project_1}. An 80\%/20\% split (random state 2026) reserved 160
observations for model development and 40 for final testing. During
hyperparameter selection, the 160 observations were split again into
120 fitting and 40 validation samples (random state 2027). The single
neural-network input feature was standardized using the fitting-set mean
and standard deviation; the same transformation was applied to the
validation observations. Following model selection, preprocessing was
refitted using all 160 training observations and applied unchanged to
the final test set. Targets were not standardized.

Two fully connected architectures were implemented in NumPy \cite{harris2020array}: one hidden
layer with 50 neurons, $[1,50,1]$, and two hidden layers of 100 neurons
each, $[1,100,100,1]$. Hidden neurons used the sigmoid function and the
regression output used an identity activation.
Weights were drawn independently from a zero-mean normal distribution
with standard deviation $1/\sqrt{n_{\mathrm{in}}}$, where
$n_{\mathrm{in}}$ is the number of inputs to the layer; biases were
initialized to zero. The same initial parameters were reused for all
optimization comparisons within each architecture.


\subsection{Backpropagation and optimization}
\label{subsec:partb_optimization}

A custom batched forward pass saved the input activations and
pre-activations of each layer. A reverse-layer implementation of
backpropagation then returned gradients for every weight matrix and
bias vector. The gradients were checked against central finite
differences, $[C(\theta+h)-C(\theta-h)]/(2h)$, with $h=10^{-6}$, on
selected randomly sampled parameter entries.

Training compared plain full-batch gradient descent, plain mini-batch
stochastic gradient descent (SGD), mini-batch RMSprop, and mini-batch
Adam. Both constant and inverse-time learning-rate schedules were
tested, giving eight configurations in total. Full-batch runs used all
120 fitting observations per update, while the stochastic runs used
mini-batches of up to 32 observations, shuffled at the beginning of each
epoch. The inverse-time schedule was
$\eta_t=\eta_0\tau/(t+\tau)$, with $\tau=10$ and $t$ counting parameter
updates. RMSprop used $\rho=0.9$; Adam used
$\beta_1=0.9$ and $\beta_2=0.999$, with $\epsilon=10^{-8}$ for both.
Plain GD and SGD were tested at
$\eta_0\in\{0.01,0.05,0.1\}$; RMSprop and Adam at
$\eta_0\in\{10^{-4},10^{-3},10^{-2}\}$. For deeper understanding of SGD, RMSprop and Adam see the Theory section in Project~1 \cite{Project_1}. Each of the 48 combinations
(two architectures, eight optimizer--schedule configurations, and three
learning rates) was trained for 300 epochs. The final validation MSE,
not the test error, determined the selected configuration for each
architecture.

In the results reported below, the selected networks were subsequently
reinitialized and fitted using all 160 training observations for
300 epochs. 


\subsection{Polynomial OLS reference}
\label{subsec:partb_ols}

The OLS baseline reproduced the best polynomial order from Project~1\cite{Project_1},
$d=14$. Its design matrix comprised $1,x,\ldots,x^{14}$. Each
non-intercept column was standardized using statistics estimated from
the same 160 training observations, and its coefficients were computed
with the Moore--Penrose pseudoinverse. OLS and the NNs were
evaluated against the identical 40 held-out targets using MSE and
$R^2$. Unlike the NNs, OLS used predefined polynomial
features, allowing a direct comparison between explicit and learned
nonlinear representations.


\subsection{Activation-function, width, and depth comparison}
\label{subsec:partc_methods}

This section extends the previous unregularized Runge regression experiment to compare three hidden-layer activations: sigmoid, ReLU, and
Leaky ReLU ($\alpha=0.01$). Four fully connected architectures were
considered: $[1,50,1]$, $[1,100,1]$, $[1,50,50,1]$, and
$[1,100,100,1]$. This $2\times2$ arrangement allows the effect of
width (50 or 100 neurons per hidden layer) to be examined separately
from depth (one or two hidden layers). All networks retained a linear
output neuron and were trained with unregularized MSE.

The same generated observations, 120 fitting samples, 40 validation
samples, and training-only input standardization were
retained, leaving the 40 test observations unused during the present
comparison. Parameters were initialized with random seed 2026 and zero
biases. Sigmoid hidden layers and linear output layers used weights
with standard deviation $1/\sqrt{n_{\mathrm{in}}}$; hidden ReLU and
Leaky ReLU layers used He initialization with standard deviation
$\sqrt{2/n_{\mathrm{in}}}$. Backpropagation was implemented by supplying the appropriate activation
functions and derivatives, and representative gradients were checked
with central finite differences.

For each of the 12 architecture--activation combinations, mini-batch
Adam was run for 300 epochs with batch size 32, constant learning rate,
and initial rates $\eta_0\in\{10^{-4},10^{-3},10^{-2}\}$. Thus, 36
training runs were compared. For each combination, the rate minimizing
\emph{validation MSE at epoch 300} was selected. Both training and
validation MSE were stored after every epoch. To examine possible
overfitting and training instability the following was recorded: the final
validation--training MSE difference, the epoch attaining the minimum
validation MSE, and that minimum value. No regularization or early
stopping was applied at this stage.

Following validation-based model selection, the architecture
and activation function with the lowest final validation MSE
were selected for held-out evaluation. The selected network
was reinitialized and trained for 300 epochs using all 160
training observations, with input scaling recalculated from
this training set. Its predictive performance was evaluated
on the same 40 test observations used for the degree-14 OLS
benchmark \cite{Project_1}.


\subsection{$L_1$ and $L_2$ regularization study}
\label{subsec:partd_methods}

The Runge dataset and the fitting/validation split described above
were retained (120 fitting, 40 validation observations); the 40
held-out test observations were not used in this experiment.
Using the ReLU activation function, four network
architectures were considered: $[1,50,1]$, $[1,100,1]$,
$[1,50,50,1]$, and $[1,100,100,1]$. Weights were initialized with
the same He scheme and random seed as in the previous step. Only weights,
not biases, were penalized. The implementation augmented the
backpropagated weight gradients with the appropriate $L_1$
subgradient or squared $L_2$ derivative (explained in Sec.~\ref{subsec:nn_regularization_theory}); both were verified against
central finite differences. The existing mini-batch Adam optimizer
was reused with batch size 32, 300 epochs, and a constant learning
rate.

The regularization strengths were
$\lambda\in\{0,10^{-8},10^{-6},10^{-4},10^{-2},10^{-1}\}$, combined with
$\eta_0\in\{10^{-3},10^{-2}\}$. The resulting grid comprised
$4\times2\times6\times2=96$ runs, including duplicate
unregularized controls ($\lambda=0$) for both penalty types.
For each configuration, the trained parameters and histories of the
unpenalized training and validation MSE were retained. Model selection
was based on \emph{validation MSE after epoch 300}, not on the
penalized training objective. The complete numerical grid is
reported in Appendix~\ref{app:partd_grid}.

For comparison with Project~1 \cite{Project_1}, polynomial Ridge with degree 10 and
$\lambda=10^{-6}$ and Lasso with degree 14 and
$\lambda=10^{-4}$ were fitted on the same 120 observations and
assessed on the same validation partition.The Lasso benchmark was fitted using the \texttt{Lasso} estimator from
\texttt{scikit-learn} rather than a custom optimization routine
\cite{Pedregosa2011}.  These polynomial configurations
were adopted from Project~1 rather than reselected using the current
validation split.


\subsection{MNIST multiclass classification}
\label{subsec:mnist_methods}

The regression framework was extended to ten-class recognition using
the full MNIST dataset \cite{lecun2010mnist}. Its canonical split comprises
$60\,000$ training and $10\,000$ test images of handwritten digits.
Each $28\times28$ grayscale image was flattened into 784 input features,
and pixel intensities were divided by 255 to map them to $[0,1]$.
The original training set was divided into $54\,000$ fitting and $6\,000$ validation images.
Integer labels were converted to ten-component one-hot vectors for the
custom implementation. The official test images were excluded from all
model selection.

The feed-forward and backpropagation code used in the previous experiments was reused,
replacing the linear regression output with ten softmax units and MSE
with the multiclass cross-entropy of
Sec.~\ref{subsec:softmax_cross_entropy}. Cross-entropy will be named CE from now on throughout the text. The classification output
gradient is given by Eq.~\eqref{eq:mnist_output_delta}; the same
hidden-layer recursion, optional weight penalties, and Adam updates
were retained \cite{Book_Ch8}. Classification gradients were checked
during development against central finite differences, as for the
regression implementation. Models were initialized with seed 2026 and
zero biases; hidden ReLU and Leaky ReLU layers used He initialization,
while sigmoid hidden layers used a standard deviation of
$1/\sqrt{n_{\mathrm{in}}}$. Unless specified otherwise, training used
mini-batch Adam, batch size 128, a constant learning rate, and 15 epochs.

A no-hidden-layer classifier $[784,10]$ served as the multinomial
logistic-regression baseline \cite{Book_Ch5}. Initially, it was compared
with $[784,100,10]$ and $[784,100,100,10]$ ReLU networks. Learning rates
$\eta_0\in\{3\times10^{-4},10^{-3},3\times10^{-3}\}$ were examined.
To investigate width, depth and activation, four architectures
($[784,50,10]$, $[784,100,10]$, $[784,50,50,10]$, and
$[784,100,100,10]$) were combined with sigmoid, ReLU and Leaky ReLU
hidden units at $\eta_0=10^{-3}$. Their epoch-15 validation accuracy
was compared, using validation CE to resolve ties.

Regularization was then investigated on the $[784,100,10]$ ReLU
network. The unregularized model was compared with
$L_1$ strengths $\lambda\in\{10^{-6},10^{-5},10^{-4}\}$ and
squared $L_2$ strengths $\lambda\in\{10^{-5},10^{-4},10^{-3}\}$.
The penalties of Sec.~\ref{subsec:nn_regularization_theory} were
applied to all weight matrices, but not biases; their sums were not
normalized by the number of parameters. Both training and validation
cross-entropies reported below exclude the penalty.

Finally, early stopping monitored the unpenalized validation
CE with patience four, retaining a copy of the weights at
its lowest observed value. The leading $L_2$ configuration was checked
with a maximum of 40 epochs. After validation-based selection, the
custom classifiers were reinitialized and fitted on all $60\,000$
training images for fixed epoch counts, without consulting the held-out
test set. A scikit-learn \texttt{MLPClassifier} with 100 ReLU neurons,
Adam and 23 fixed epochs was also trained on the $54\,000$ fitting
images as an external \emph{validation} reference \cite{Pedregosa2011}.
It used no weight penalty or validation-selected checkpoint and hence
was not a fully matched comparison. The selected custom classifier was
finally evaluated on the $10\,000$ test images using accuracy,
CE and a confusion matrix.


% ============================================================
\section{Results and Discussion}


\subsection{Gradient verification and optimizer sensitivity}
\label{subsec:partb_results_optim}

Numerical differentiation confirmed the backpropagation gradients at
the sampled entries. For the $[1,50,1]$ network, 16 comparisons gave a
maximum absolute difference of $1.339\times10^{-10}$; for the
$[1,100,100,1]$ network, 26 comparisons gave
$2.310\times10^{-11}$. No comparison failed the prescribed numerical
tolerance. 

Figures~\ref{fig:partb_learning_rates} and
\ref{fig:partb_validation_curves} show the sensitivity to optimizer and
learning rate. Plain full-batch GD and
mini-batch SGD show large early
fluctuations in Fig~\ref{fig:partb_validation_curves} (purple and blue curves). RMSprop (red curve in Fig.~\ref{fig:partb_validation_curves}) reached lower
errors but produced more irregular validation curves, especially for [1, 100, 100, 1] architecture.
Adam (green curve in Fig.~\ref{fig:partb_validation_curves}) with a constant learning rate attained the smallest final
validation error for both architectures. Inverse-time scheduling often
reduced the step size before the networks had fitted the narrow central
peak of Runge's function, leaving relatively high validation errors;
this behaviour is specific to the tested decay scale and training
budget rather than an intrinsic limitation of decreasing schedules.

\begin{figure*}[t]
 \centering
 \includegraphics[width=0.49\textwidth]{Figures/NN_learning_rates_1-50-1.pdf}\hfill
 \includegraphics[width=0.49\textwidth]{Figures/NN_learning_rates_1-100-100-1.pdf}
 \caption{Final validation MSE as a function of initial learning rate
 for the $[1,50,1]$ (left) and $[1,100,100,1]$ (right) networks.
 Each curve represents one optimizer--learning-rate-schedule
 configuration after 300 epochs.}
 \label{fig:partb_learning_rates}
\end{figure*}

\begin{figure*}[t]
 \centering
 \includegraphics[width=0.49\textwidth]{Figures/NN_optimizer_validation_1-50-1.pdf}\hfill
 \includegraphics[width=0.49\textwidth]{Figures/NN_optimizer_validation_1-100-100-1.pdf}
 \caption{Validation MSE over 300 epochs for the best tested initial
 learning rate within each optimizer--schedule configuration. Solid
 curves correspond to constant learning rates and dashed curves to
 inverse-time decay. The vertical axis is logarithmic.}
 \label{fig:partb_validation_curves}
\end{figure*}

The selected configuration was constant-rate mini-batch Adam with
$\eta_0=0.01$ for both models. The final validation MSE was 0.011731
for $[1,50,1]$ and 0.011246 for $[1,100,100,1]$. These are
validation-set model-selection results, not measurements on the held-out
test set.


\subsection{Comparison with polynomial OLS}
\label{subsec:partb_results_ols}

Table~\ref{tab:partb_model_comparison} summarizes the final evaluation
on the shared 80\%/20\% split. For the reported 300-epoch final
fitting, both network architectures performed slightly worse than the best OLS case.
The differences are modest and concern one realization with only 40
test observations.

\begin{table}[h]
 \centering
 \caption{Training and held-out test performance for the degree-14 OLS
 reference and two feed-forward networks. The
neural-network entries correspond to 300 final-training epochs.}
 \label{tab:partb_model_comparison}
 \begin{tabular}{lccc}
 \hline\hline
 Model & Train MSE & Test MSE & Test $R^2$ \\
 \hline
 OLS ($d=14$) & 0.009278 & 0.010240 & 0.867169 \\
 NN $[1,50,1]$ & 0.009931 & 0.011101 & 0.856006 \\
 NN $[1,100,100,1]$ & 0.010264 & 0.011555 & 0.850116 \\
 \hline\hline
 \end{tabular}
\end{table}

Figure~\ref{fig:partb_runge_fits} compares the fitted curves with the
exact Runge function and the noisy observations. All three models
reproduce the narrow peak near $x=0$ and its decay toward the interval
edges. The deeper network, referred to from now on as Network~2, corresponding to architecture [1, 100, 100, 1], slightly overestimates the peak. The simpler network, Network 1, corresponding to architecture [1, 50, 1], replicates the peak the best but is farther from the exact Runge function than OLS in the lateral slopes. The degree-14 polynomial fit is also accurate in the central
region but displays more pronounced boundary oscillations, especially
near $x=\pm1$, where polynomial extrapolation near the edges is
sensitive to the fitted coefficients. 

\begin{figure}[t]
 \centering
 \includegraphics[width=0.49\textwidth]{Figures/Runge_OLS_d14_vs_NN.pdf}
 \caption{Exact Runge function, noisy training and test observations,
 degree-14 OLS, and the two trained NNs, all evaluated on
 the same dense input grid. The displayed networks correspond to the
 300-epoch final fits reported in Table~\ref{tab:partb_model_comparison}.}
 \label{fig:partb_runge_fits}
\end{figure}

OLS uses only 15 fitted polynomial coefficients, whereas the two neural
networks have 151 and 10\,401 trainable weights and biases, respectively.

The slightly better performance of OLS can be explained by the
relatively simple, smooth, one-dimensional nature of Runge's
function, which can be accurately approximated by a degree-14
polynomial using only 15 coefficients. In contrast, the neural
networks contain substantially more trainable parameters relative
to the 160 training observations, particularly Network 2, which
may increase their susceptibility to fitting noise and overfitting.
However, the higher training MSE of Network 2 compared with OLS
suggests that optimization effects may also contribute to its
poorer test performance. Thus, for this particular regression
problem, the simpler polynomial model provides a favorable
balance between model complexity and generalization.


\subsection{Activation functions, network depth, and width}
\label{subsec:partc_results}

Table~\ref{tab:partc_activation_mse} and
Fig.~\ref{fig:partc_architecture} summarize the final validation
MSE for all twelve architecture--activation combinations. The selected
initial learning rate was $\eta_0=0.01$ in every case, the largest of
the three rates tested. The lowest \emph{final} validation MSE,
$0.010980$, was obtained with ReLU and two hidden layers of 100 neurons.
Nevertheless, the $[1,100,1]$ sigmoid network achieved $0.011004$,
a difference of only $2.4\times10^{-5}$. Thus, the reported ranking
is sensitive to small performance differences and should not be
interpreted as evidence of general superiority from one split and
initialization.

\begin{table}[h]
 \centering
 \caption{Validation MSE at epoch 300 for each architecture and
 hidden-layer activation. All selected learning rates were $0.01$.
 The smallest entry in each row is bold.}
 \label{tab:partc_activation_mse}
 \begin{tabular}{lccc}
 \hline\hline
 Architecture & Sigmoid & ReLU & Leaky ReLU \\
 \hline
 $[1,50,1]$ & 0.011731 & 0.011501 & \textbf{0.011477} \\
 $[1,100,1]$ & \textbf{0.011004} & 0.011267 & 0.011210 \\
 $[1,50,50,1]$ & \textbf{0.011178} & 0.011346 & 0.011404 \\
 $[1,100,100,1]$ & 0.011246 & \textbf{0.010980} & 0.011055 \\
 \hline\hline
 \end{tabular}
\end{table}

\begin{figure}[t]
 \centering
 \includegraphics[width=0.49\textwidth]{Figures/NN_activation_architecture_comparison.pdf}
 \caption{Final validation MSE after 300 epochs for sigmoid, ReLU,
 and Leaky ReLU hidden layers in networks of different depths and
 widths. All twelve configurations selected an Adam learning rate
 of $\eta_0=0.01$ from the tested grid. }
 \label{fig:partc_architecture}
\end{figure}

Increasing width or depth did not consistently reduce validation
error. Sigmoid performed best within the
$[1,100,1]$ and $[1,50,50,1]$ architectures, whereas Leaky ReLU
performed best within $[1,50,1]$ and ReLU within $[1,100,100,1]$.
The influence of model capacity therefore depends on the activation
and its optimization trajectory; more trainable parameters alone did
not ensure better generalization.

Figure~\ref{fig:partc_learning_curves} shows the training and
validation histories, restricted to epochs 30--300 to make the later
behavior visible. The sigmoid curves display a comparatively smooth
initial decrease over the displayed interval. ReLU and Leaky ReLU
usually reach the low-error regime sooner but exhibit more frequent
spikes and fluctuations, particularly in the wider and deeper
networks. The first 30 epochs are not shown, so the plots
do not establish which activation had the sharpest decrease at the
very start of training.

\begin{figure*}[t]
 \centering
 \includegraphics[width=0.48\textwidth]{Figures/NN_C_learning_curves_1-50-1.pdf}\hfill
 \includegraphics[width=0.48\textwidth]{Figures/NN_C_learning_curves_1-100-1.pdf}\\[1ex]
 \includegraphics[width=0.48\textwidth]{Figures/NN_C_learning_curves_1-50-50-1.pdf}\hfill
 \includegraphics[width=0.48\textwidth]{Figures/NN_C_learning_curves_1-100-100-1.pdf}
 \caption{Training (dashed) and validation (solid) MSE histories for
 one-hidden-layer networks with 50 and 100 neurons (top row) and
 two-hidden-layer networks with 50 and 100 neurons per layer (bottom
 row). Each activation uses its best tested initial learning rate;
 the displayed interval is epochs 30--300.}
 \label{fig:partc_learning_curves}
\end{figure*}

The final generalization gaps and best observed validation epochs are
reported in Table~\ref{tab:partc_diagnostics}. Note that some validation trajectories attained lower errors well before the
last epoch. Overall, the training--validation gaps and intermittent deterioration in validation MSE provide \emph{indications}, but not a quantitative bias--variance decomposition, of possible overfitting. 

\begin{table*}[t]
 \centering
 \caption{Generalization and optimization diagnostics for the learning
 rate selected within each architecture--activation combination. The
 gap is final validation MSE minus final training MSE. The minimum
 validation MSE can occur before epoch 300.}
 \label{tab:partc_diagnostics}
 \begin{ruledtabular}
 \begin{tabular}{llccc} %{llrrr}
 Architecture & Activation & Final gap & Best epoch & Min. validation MSE \\
 \hline
 $[1,50,1]$ & Sigmoid & 0.001350 & 277 & 0.010642 \\
                & ReLU & 0.002482 & 296 & 0.010725 \\
                & Leaky ReLU & 0.002391 & 296 & 0.010669 \\
 $[1,100,1]$ & Sigmoid & 0.001868 & 298 & 0.010852 \\
                & ReLU & 0.002435 & 100 & 0.010994 \\
                & Leaky ReLU & 0.002680 & 120 & 0.010930 \\
 $[1,50,50,1]$ & Sigmoid & 0.001716 & 220 & 0.010450 \\
                & ReLU & 0.003728 & 56 & 0.009360 \\
                & Leaky ReLU & 0.003780 & 56 & 0.009856 \\
 $[1,100,100,1]$ & Sigmoid & 0.001357 & 282 & 0.010279 \\
                & ReLU & 0.003054 & 156 & 0.010737 \\
                & Leaky ReLU & 0.003162 & 80 & 0.009924 \\
 \end{tabular}
 \end{ruledtabular}
\end{table*}

Finally, the best configuration selected from the validation
experiments was the ReLU network with architecture
$[1,100,100,1]$ and learning rate $\eta_0=0.01$.
After retraining on all 160 training observations (leaving 40 for testing) for 300
epochs, this model achieved a test MSE of $0.008492$ and
$R^2=0.889851$, outperforming both sigmoid networks and the degree-14 OLS benchmark from the previous comparison (see Table~\ref{tab:partb_model_comparison}).

Figure~\ref{fig:partc_runge_fits} compares the selected ReLU
network with OLS and the exact Runge function. The ReLU
network shows some visible changes in
slope due to the piecewise-linear nature of ReLU.
The polynomial OLS fit is comparatively smooth in the
central region but exhibits larger deviations near the
interval boundaries. 

\begin{figure}[t]
    \centering
    \includegraphics[width=\columnwidth]
    {Figures/Runge_OLS_d14_vs_NN_2.pdf}
    \caption{Exact Runge function,
    noisy training and test observations, degree-14
    polynomial OLS, and the selected ReLU NN
    with architecture $[1,100,100,1]$. Both models were
    fitted using the same noisy 160 training observations and
    evaluated on the same 40 held-out test observations.}
    \label{fig:partc_runge_fits}
\end{figure}


\subsection{Effect of $L_1$ and $L_2$ regularization}
\label{subsec:partd_results}

Table~\ref{tab:partd_architecture_results} reports the lowest
\emph{final validation MSE} for each architecture and penalty type ($L_1$ or $L_2$),
selected over the tested learning rates and regularization strengths (see Appendix~\ref{app:partd_grid} for more results).
Both penalties attained their lowest overall validation MSE with the
two-hidden-layer network $[1,100,100,1]$. However, the preferred
regularization strength depended on the architecture and penalty:
stronger penalization was not uniformly beneficial. In particular,
for the one-hidden-layer network $[1,100,1]$, the best $L_2$
configuration had $\lambda=0$.

\begin{table}[h]
\centering
\caption{Best validation MSE at epoch 300 for each ReLU network and
penalty, selected over the tested $\lambda$ and learning-rate values.}
\label{tab:partd_architecture_results}
\begin{ruledtabular}
\begin{tabular}{llccc}
Architecture & Penalty & $\lambda$ & $\eta_0$ & Val. MSE \\
$[1,50,1]$ & $L_1$ & $10^{-4}$ & 0.01 & 0.010932 \\
$[1,50,1]$ & $L_2$ & $10^{-4}$ & 0.01 & 0.011104 \\
$[1,100,1]$ & $L_1$ & $10^{-8}$ & 0.01 & 0.011264 \\
$[1,100,1]$ & $L_2$ & $0$ & 0.01 & 0.011267 \\
$[1,50,50,1]$ & $L_1$ & $10^{-8}$ & 0.01 & 0.011293 \\
$[1,50,50,1]$ & $L_2$ & $10^{-2}$ & 0.001 & 0.010407 \\
$[1,100,100,1]$ & $L_1$ & $10^{-4}$ & 0.01 & 0.010556 \\
$[1,100,100,1]$ & $L_2$ & $10^{-2}$ & 0.001 & 0.010249 \\
\end{tabular}
\end{ruledtabular}
\end{table}

Table~\ref{tab:partd_reference_results} compares the selected
network configurations with the unregularized ReLU reference and the
Project~1 polynomial Ridge and Lasso models on the same validation
observations. For squared $L_2$, $\lambda=10^{-2}$ and
$\eta_0=0.001$ gave the lowest observed validation MSE,
$0.010249$ (about $6.7\%$ below the unregularized reference).
The historical
Ridge and Lasso configurations yielded validation MSEs of
$0.012406$ and $0.012875$, respectively.

\begin{table}[h]
\centering
\caption{Best observed regularized and unregularized ReLU NNs,
and polynomial regression references (Lasso and Ridge) on the fitting/validation split.
The NN use architecture $[1,100,100,1]$.}
\label{tab:partd_reference_results}
\begin{ruledtabular}
\begin{tabular}{lccc}
Model & $\lambda$ & Train MSE & Val. MSE \\
Unreg. NN & $0$ & 0.007927 & 0.010980 \\
$L_1$ NN & $10^{-4}$ & 0.008143 & 0.010556 \\
$L_2$ NN & $10^{-2}$ & 0.009133 & \textbf{0.010249} \\
$L_1$ Lasso ($d=14$) & $10^{-4}$ & --- & 0.012875 \\
$L_2$ Ridge ($d=10$) & $10^{-6}$ & --- & 0.012406 \\
\end{tabular}
\end{ruledtabular}
\end{table}

Figure~\ref{fig:partd_regularization_curves} shows the fitting and
validation trajectories of the two selected regularized networks
from epochs 30--300. The $L_1$ network, trained with the larger
learning rate, exhibits appreciable fluctuations, especially in
validation MSE. The $L_2$ network shows a more gradual reduction
followed by comparatively stable errors. Their final
validation--training MSE differences were $0.002413$ and $0.001116$
for $L_1$ and $L_2$, respectively, versus $0.003053$ for the
unregularized reference. These differences are consistent with a
reduced generalization gap, but they do not by themselves establish
that the networks are free of overfitting.

\begin{figure*}[t]
 \centering
 \includegraphics[width=0.85\textwidth]{Figures/NN_D_regularization_learning_curves.pdf}
 \caption{Unpenalized training and validation MSE from epochs
 30--300 for the selected $[1,100,100,1]$ ReLU networks with
 $L_1$ ($\lambda=10^{-4}$, $\eta_0=0.01$) and squared $L_2$
 ($\lambda=10^{-2}$, $\eta_0=0.001$) penalties.}
 \label{fig:partd_regularization_curves}
\end{figure*}

\noindent\textit{Final held-out evaluation.}

Following validation-based hyperparameter selection,
the best unregularized, $L_1$-regularized, and
$L_2$-regularized ReLU NNs were reinitialized
and trained for 300 epochs using all 160 training
observations. All three selected networks had
architecture $[1,100,100,1]$. Their predictions
were evaluated on the same 40 held-out test
observations as the polynomial regression references.

Table~\ref{tab:partd_final_comparison} summarizes the
final performance. Although $L_2$ regularization
achieved the lowest validation MSE during model
selection, the unregularized network obtained the
lowest test MSE ($0.008492$), followed closely by the
$L_1$ network ($0.008555$). The $L_2$ network
achieved a higher test MSE ($0.009770$), despite
its superior validation performance. Nevertheless,
all three neural networks outperformed the
polynomial OLS, Ridge, and Lasso references on
this particular test split.

\begin{table*}[t]
\centering
\caption{Final training and held-out test performance
of the polynomial regression references and selected
ReLU NNs. All neural networks use
architecture $[1,100,100,1]$ and were retrained
for 300 epochs on the complete training set.}
\label{tab:partd_final_comparison}
\begin{ruledtabular}
\begin{tabular}{lccc}
Model & Train MSE & Test MSE & Test $R^2$ \\
\hline
OLS ($d=14$)       & 0.009278 & 0.010240 & 0.867169 \\
Ridge ($d=10$)     & 0.010720 & 0.011360 & 0.852644 \\
Lasso ($d=14$)     & 0.011544 & 0.012094 & 0.843116 \\
Unregularized NN   & 0.008884 & \textbf{0.008492} & \textbf{0.889851} \\
$L_1$ NN           & 0.009026 & 0.008555 & 0.889033 \\
$L_2$ NN           & 0.009698 & 0.009770 & 0.873262 \\
\end{tabular}
\end{ruledtabular}
\end{table*}

Figure~\ref{fig:partd_final_fits} compares the fitted
functions obtained with $L_1$ and $L_2$
regularization against their corresponding
polynomial regression methods.

The $L_1$ neural network reproduces the narrow
central peak of Runge's function more accurately
than Lasso, which underestimates its maximum
and exhibits deviations near the interval boundary.
For $L_2$, both the neural network and Ridge
underestimate the central peak, although the
neural network achieves a lower overall test MSE.
The ReLU networks also exhibit local changes in
slope due to their piecewise-linear activations.

\begin{figure*}[t]
    \centering
    \includegraphics[width=0.95\textwidth]
    {Figures/NN_D_final_L1_L2_vs_Ridge_Lasso.pdf}
    \caption{Comparison of the exact Runge function,
    noisy training and test observations, and fitted
    regression models. Left: $L_1$-regularized ReLU
    NN ($\lambda=10^{-4}$) versus degree-14
    Lasso. Right: $L_2$-regularized ReLU NN
    ($\lambda=10^{-2}$) versus degree-10 Ridge.
    The NNs use architecture
    $[1,100,100,1]$. All models were fitted on the
    same noisy 160 training observations and evaluated
    on the same 40 held-out test observations.}
    \label{fig:partd_final_fits}
\end{figure*}

These results show that the improvements obtained
through regularization during validation were not
fully reproduced on the held-out test set.
In particular, the test MSE of the $L_1$ network
was only slightly higher than that of the
unregularized model, whereas $L_2$ produced a
larger deterioration. The different ranking may
reflect sensitivity to the small validation and
test subsets, the retraining procedure, and
stochastic optimization. Thus, although
regularization can improve validation performance,
it did not provide a consistent advantage over
the unregularized ReLU network in the present
experiment.


\subsection{MNIST multiclass classification}
\label{subsec:mnist_results}

\noindent\textit{Baseline and learning-rate sensitivity.}


The initial comparison (Table~\ref{tab:mnist_baseline}) shows that
adding a hidden layer substantially improved classification over the
multinomial logistic-regression baseline. After 15 epochs at
$\eta_0=10^{-3}$, validation accuracy increased from $92.75\%$ for the
baseline to $97.93\%$ for the one-hidden-layer ReLU network. Adding a
second hidden layer did not help: the $[784,100,100,10]$ network
achieved $97.67\%$ and had a higher validation CE despite a
lower training CE ($0.01109$ versus $0.01918$). Its
validation CE reached its minimum at epoch 5 but rose to
$0.11109$ by epoch 15, providing evidence of overfitting under these
settings. Figure~\ref{fig:mnist_learning_curves} illustrates the
training and validation trajectories.

\begin{table}[h]
\centering
\caption{MNIST fitting/validation comparison after 15 epochs with
$\eta_0=10^{-3}$, Adam and batch size 128. All hidden layers use ReLU.}
\label{tab:mnist_baseline}
\begin{ruledtabular}
\begin{tabular}{lcc}
Model & Val. CE & Val. accuracy \\
Logistic $[784,10]$ & 0.264513 & $92.75\%$ \\
NN $[784,100,10]$ & \textbf{0.080613} & $\mathbf{97.93\%}$ \\
NN $[784,100,100,10]$ & 0.111089 & $97.67\%$ \\
\end{tabular}
\end{ruledtabular}
\end{table}

\begin{figure*}[t]
\centering
\includegraphics[width=0.93\textwidth]{Figures/MNIST_E_learning_curves.pdf}
\caption{Training (dashed) and validation (solid) CE and
accuracy for multinomial logistic regression and two ReLU networks.
All models use Adam with $\eta_0=10^{-3}$ and 15 epochs; the initial
state (epoch 0) is omitted to make the learning trajectories clearer. NN 1 corresponds to architecture [784, 100, 10] and NN 2 corresponds to architecture [784, 100, 100, 10].}
\label{fig:mnist_learning_curves}
\end{figure*}

The learning-rate study found $\eta_0=10^{-3}$ to be the best tested
rate for both hidden-layer networks in terms of final validation
accuracy. For $[784,100,10]$, increasing $\eta_0$ to $3\times10^{-3}$
reduced training CE but worsened validation CE
from $0.08061$ to $0.11786$ and validation accuracy from $97.93\%$ to
$97.72\%$. For logistic regression, $\eta_0=3\times10^{-3}$ gave a
slightly higher validation accuracy ($92.80\%$) than $\eta_0=10^{-3}$
($92.75\%$), although its CE was higher ($0.27188$ versus
$0.26451$). This illustrates that accuracy and CE need not
rank configurations identically.

\noindent\textit{Activation functions, depth and width.}


Figure~\ref{fig:mnist_architecture} summarizes the 12
architecture--activation combinations trained at $\eta_0=10^{-3}$.
The one-hidden-layer $[784,100,10]$ (1x100) ReLU network was best by both
final validation accuracy ($97.93\%$) and CE ($0.08061$).
Its Leaky ReLU counterpart was close ($97.87\%$, CE $0.08298$).
For two 100-neuron hidden layers (2x100), Leaky ReLU gave the highest accuracy
($97.80\%$), whereas sigmoid gave a lower CE ($0.08210$).
The effect of increased depth therefore depended on the activation,
and more layers or neurons did not consistently improve validation
performance. All comparisons used one data split and initialization,
so small differences should not be interpreted as universal rankings.

\begin{figure*}[t]
\centering
\includegraphics[width=0.89\textwidth]{Figures/MNIST_E_activation_architecture_heatmaps.pdf}
\caption{MNIST validation accuracy (left) and CE (right)
after 15 epochs for four network architectures and three hidden-layer
activations. All runs use mini-batch Adam, batch size 128 and
$\eta_0=10^{-3}$. Architecture labels indicate the number and width of
hidden layers.}
\label{fig:mnist_architecture}
\end{figure*}

\noindent\textit{Regularization and early stopping.}

At a fixed budget of 15 epochs, weak $L_2$ regularization
($\lambda=10^{-5}$) produced the highest validation accuracy in the
regularization search ($97.98\%$); a stronger $L_2$ penalty
($\lambda=10^{-4}$) achieved the lowest CE ($0.07564$),
with $97.92\%$ accuracy (Table~\ref{tab:mnist_regularization}).
The $L_1$ model with $\lambda=10^{-5}$ attained a similar CE of
$0.07579$, while stronger penalties worsened performance. The gain
in accuracy from $97.93\%$ to $97.98\%$ corresponds to only three
additional correct images among the 6000 validation observations.

\begin{table}[h]
\centering
\caption{Validation results for the $[784,100,10]$ (1x100) ReLU network
after 15 epochs at $\eta_0=10^{-3}$. CE excludes the regularization
penalty.}
\label{tab:mnist_regularization}
\begin{ruledtabular}
\begin{tabular}{lccc}
Penalty & $\lambda$ & Val. CE & Val. acc. \\
None & 0 & 0.08061 & $97.93\%$ \\
$L_1$ & $10^{-6}$ & 0.08008 & $97.93\%$ \\
$L_1$ & $10^{-5}$ & 0.07579 & $97.87\%$ \\
$L_1$ & $10^{-4}$ & 0.08689 & $97.40\%$ \\
$L_2$ & $10^{-5}$ & 0.07902 & $\mathbf{97.98\%}$ \\
$L_2$ & $10^{-4}$ & \textbf{0.07564} & $97.92\%$ \\
$L_2$ & $10^{-3}$ & 0.11900 & $97.00\%$ \\
\end{tabular}
\end{ruledtabular}
\end{table}

Allowing early stopping with patience four changed the comparison:
the unregularized model reached a minimum validation CE of $0.07802$
at epoch 13 ($97.83\%$ accuracy), the $L_2$ model with
$\lambda=10^{-5}$ reached $0.07751$ at epoch 16 ($97.90\%$), and
$\lambda=10^{-4}$ reached $0.06784$ at epoch 23 ($98.07\%$).
When the maximum training budget for the last configuration was
extended to 40 epochs, training stopped at epoch 27 and selected the
same epoch-23 checkpoint. These results favored the latter model for
final evaluation; the checkpoint weights were verified by recomputing
its validation CE independently.

As a library reference, a scikit-learn \texttt{MLPClassifier} with
architecture $[784,100,10]$, ReLU, Adam and 23 fixed epochs achieved
$97.80\%$ validation accuracy and CE $0.09142$, compared with
$98.07\%$ and $0.06784$ for the custom selected checkpoint
\cite{Pedregosa2011}. However, the scikit-learn model was
unregularized and did not use a validation-selected checkpoint.
Different initialization and optimization trajectories also prevent
attributing this gap to an intrinsic advantage of the custom code.
These are \emph{validation} results; no scikit-learn test result was
computed in the final evaluation.

\noindent\textit{Held-out test performance and error structure.}

After hyperparameters and epoch counts were fixed, the custom models
were reinitialized and trained on all $60\,000$ MNIST training images.
Table~\ref{tab:mnist_test} reports their performance on the untouched
$10\,000$-image official test set. The $L_2$-regularized ReLU network
achieved $97.90\%$ test accuracy and CE $0.064215$, compared with
$97.87\%$ and CE $0.071666$ for its unregularized counterpart. Both
networks substantially outperformed multinomial logistic regression. Nevertheless, the accuracy difference
between the two neural networks is only $0.03$ percentage points,
equivalent to three images; a single split and seed do not establish a
statistically reliable advantage of $L_2$ regularization. The lower
test CE and smaller training--test accuracy gap are
consistent with improved generalization in this particular run.

\begin{table*}[t]
\centering
\caption{Final MNIST training and held-out test results after refitting
on all $60\,000$ training images. The official $10\,000$ test images
were not used for hyperparameter selection. The selected $L_2$ NN uses
$\lambda=10^{-4}$ and $\eta_0=10^{-3}$.}
\label{tab:mnist_test}
\begin{ruledtabular}
\begin{tabular}{lccccc}
Model & Epochs & Train CE & Test CE & Train accuracy & Test accuracy \\
Logistic regression $[784,10]$ & 15 & 0.236807 & 0.274847 & $93.48\%$ & $92.45\%$ \\
Unregularized NN $[784,100,10]$ & 15 & 0.014306 & 0.071666 & $99.74\%$ & $97.87\%$ \\
$L_2$ NN $[784,100,10]$ & 23 & 0.023611 & \textbf{0.064215} & $99.56\%$ & $\mathbf{97.90\%}$ \\
\end{tabular}
\end{ruledtabular}
\end{table*}

The confusion matrix for the validation-selected $L_2$ network
(Fig.~\ref{fig:mnist_confusion}) contains $9790$ correctly classified
images and $210$ errors. The most frequent off-diagonal entries were
true $9$ predicted as $4$ (11 cases), true $4$ predicted as $9$
(10 cases), and true $7$ predicted as $2$ (9 cases). These indicate
specific digit pairs with overlapping visual features. Class $9$ had
the lowest recall ($95.94\%$) and $F_1$ score ($0.9690$); the
macro-averaged $F_1$ was $0.9789$. Thus, even with high overall
accuracy, recognition difficulty is not uniform across digits.

\begin{figure}[t]
\centering
\includegraphics[width=0.42\textwidth]{Figures/MNIST_E_confusion_matrix.pdf}
\caption{Confusion matrix of the validation-selected, $L_2$-regularized
$[784,100,10]$ ReLU network on the 10\,000-image official MNIST test
set, after refitting on all training images. Rows give true digits and
columns predicted digits; off-diagonal entries indicate errors.}
\label{fig:mnist_confusion}
\end{figure}


% ============================================================

\newpage
\vspace{10cm}
\section{Conclusions and Perspectives}


The implementation developed in this project successfully handled
both continuous regression and multiclass classification within a
common feed-forward and backpropagation framework. By adapting the
output activation and loss function, the same numerical routines
could be applied to substantially different learning problems.
Finite-difference gradient checks supported the correctness of
the implementation, while the experiments demonstrated the
sensitivity of neural networks to optimization settings,
architecture, and regularization.

For the Runge regression problem, the unregularized ReLU network
achieved the lowest observed test MSE ($0.008492$). Nevertheless,
the degree-14 polynomial OLS model remained competitive
($0.010240$) while requiring substantially fewer parameters
and offering a simpler, more interpretable representation.
Ridge and Lasso provided additional control over model
complexity but did not outperform these models in the present
experiment. In contrast, the MNIST classification task
demonstrated a clearer benefit from nonlinear hidden layers:
the selected $L_2$-regularized network achieved $97.90\%$
test accuracy, compared with $92.45\%$ for multinomial
logistic regression. However, its improvement over the
unregularized neural network was only $0.03$ percentage points.

These contrasting results illustrate that the flexibility of neural networks
does not guarantee superior performance, and model selection
should depend on the structure and complexity of the data
\cite{Book_Ch8}. For a relatively simple, low-dimensional
problem such as Runge regression, polynomial models can provide
an effective alternative at considerably lower complexity.
For a larger and more complex classification problem such as
MNIST, nonlinear neural networks offer substantial predictive
advantages. Nevertheless, our fully connected networks operated
on flattened images and did not explicitly exploit their
spatial structure.

The experiments also highlight limitations associated with
neural networks, including computational cost, sensitivity
to hyperparameters and initialization, and limited
interpretability. The results are restricted to individual
random seeds and fixed data partitions, so small performance
differences should not be regarded as conclusive. Future work
could evaluate stability across multiple seeds, compare
computational costs and probability calibration, and investigate
architectures that explicitly exploit spatial information,
such as convolutional neural networks. Applying the developed
code to other datasets would further test its generality.

Overall, the main lesson of this project is not that one
algorithm is universally superior, but that predictive
performance must be balanced against model complexity,
interpretability, computational cost, and the available data.

% ============================================================
%TC:ignore

\appendix


\section{Complete $L_1$ and $L_2$ hyperparameter results}
\label{app:partd_grid}

Tables~\ref{tab:partd_l1_full} and \ref{tab:partd_l2_full} list all 96 training runs used for the results in Sec.~\ref{subsec:partd_results}.
Values are unpenalized MSE after 300 training epochs on the
120-observation fitting and 40-observation validation sets.
The $\lambda=0$ cases appear under both penalties as an internal
reproducibility control.
\begin{table*}[p]
\centering
\caption{Complete $L_1$ grid: unpenalized training and validation MSE
at epoch 300, for four NN architectures, six $\lambda$ values,
and two learning rates.}
\label{tab:partd_l1_full}
\scriptsize
\begin{ruledtabular}
\begin{tabular}{lccrr}
Architecture & $\lambda$ & $\eta_0$ & Train MSE & Val. MSE \\
$[1,50,1]$ & $0$ & 0.001 & 0.013271 & 0.013980 \\
$[1,50,1]$ & $0$ & 0.01 & 0.009019 & 0.011501 \\
$[1,50,1]$ & $10^{-8}$ & 0.001 & 0.013271 & 0.013980 \\
$[1,50,1]$ & $10^{-8}$ & 0.01 & 0.009019 & 0.011501 \\
$[1,50,1]$ & $10^{-6}$ & 0.001 & 0.013272 & 0.013980 \\
$[1,50,1]$ & $10^{-6}$ & 0.01 & 0.009019 & 0.011518 \\
$[1,50,1]$ & $10^{-4}$ & 0.001 & 0.013296 & 0.013993 \\
$[1,50,1]$ & $10^{-4}$ & 0.01 & 0.008754 & 0.010932 \\
$[1,50,1]$ & $10^{-2}$ & 0.001 & 0.025158 & 0.025430 \\
$[1,50,1]$ & $10^{-2}$ & 0.01 & 0.029405 & 0.024923 \\
$[1,50,1]$ & $10^{-1}$ & 0.001 & 0.049263 & 0.042792 \\
$[1,50,1]$ & $10^{-1}$ & 0.01 & 0.099490 & 0.088418 \\
$[1,100,1]$ & $0$ & 0.001 & 0.012216 & 0.013462 \\
$[1,100,1]$ & $0$ & 0.01 & 0.008833 & 0.011267 \\
$[1,100,1]$ & $10^{-8}$ & 0.001 & 0.012216 & 0.013462 \\
$[1,100,1]$ & $10^{-8}$ & 0.01 & 0.008832 & 0.011264 \\
$[1,100,1]$ & $10^{-6}$ & 0.001 & 0.012217 & 0.013463 \\
$[1,100,1]$ & $10^{-6}$ & 0.01 & 0.008841 & 0.011269 \\
$[1,100,1]$ & $10^{-4}$ & 0.001 & 0.012429 & 0.013606 \\
$[1,100,1]$ & $10^{-4}$ & 0.01 & 0.008961 & 0.012547 \\
$[1,100,1]$ & $10^{-2}$ & 0.001 & 0.022570 & 0.020599 \\
$[1,100,1]$ & $10^{-2}$ & 0.01 & 0.028982 & 0.024667 \\
$[1,100,1]$ & $10^{-1}$ & 0.001 & 0.040916 & 0.035704 \\
$[1,100,1]$ & $10^{-1}$ & 0.01 & 0.099489 & 0.088449 \\
$[1,50,50,1]$ & $0$ & 0.001 & 0.008226 & 0.012321 \\
$[1,50,50,1]$ & $0$ & 0.01 & 0.007618 & 0.011346 \\
$[1,50,50,1]$ & $10^{-8}$ & 0.001 & 0.008312 & 0.011719 \\
$[1,50,50,1]$ & $10^{-8}$ & 0.01 & 0.007716 & 0.011293 \\
$[1,50,50,1]$ & $10^{-6}$ & 0.001 & 0.008294 & 0.011674 \\
$[1,50,50,1]$ & $10^{-6}$ & 0.01 & 0.007884 & 0.011435 \\
$[1,50,50,1]$ & $10^{-4}$ & 0.001 & 0.008537 & 0.011722 \\
$[1,50,50,1]$ & $10^{-4}$ & 0.01 & 0.008105 & 0.011480 \\
$[1,50,50,1]$ & $10^{-2}$ & 0.001 & 0.011014 & 0.011655 \\
$[1,50,50,1]$ & $10^{-2}$ & 0.01 & 0.012606 & 0.013254 \\
$[1,50,50,1]$ & $10^{-1}$ & 0.001 & 0.099489 & 0.088438 \\
$[1,50,50,1]$ & $10^{-1}$ & 0.01 & 0.099492 & 0.088401 \\
$[1,100,100,1]$ & $0$ & 0.001 & 0.007765 & 0.011603 \\
$[1,100,100,1]$ & $0$ & 0.01 & 0.007927 & 0.010980 \\
$[1,100,100,1]$ & $10^{-8}$ & 0.001 & 0.007837 & 0.010828 \\
$[1,100,100,1]$ & $10^{-8}$ & 0.01 & 0.007602 & 0.011737 \\
$[1,100,100,1]$ & $10^{-6}$ & 0.001 & 0.008371 & 0.012311 \\
$[1,100,100,1]$ & $10^{-6}$ & 0.01 & 0.007883 & 0.010897 \\
$[1,100,100,1]$ & $10^{-4}$ & 0.001 & 0.008293 & 0.011063 \\
$[1,100,100,1]$ & $10^{-4}$ & 0.01 & 0.008143 & 0.010556 \\
$[1,100,100,1]$ & $10^{-2}$ & 0.001 & 0.012197 & 0.011967 \\
$[1,100,100,1]$ & $10^{-2}$ & 0.01 & 0.012551 & 0.013209 \\
$[1,100,100,1]$ & $10^{-1}$ & 0.001 & 0.099489 & 0.088438 \\
$[1,100,100,1]$ & $10^{-1}$ & 0.01 & 0.099491 & 0.088412 \\
\end{tabular}
\end{ruledtabular}
\end{table*}
\begin{table*}[p]
\centering
\caption{Complete $L_2$ grid: unpenalized training and validation MSE
at epoch 300, for four NN architectures, six $\lambda$ values,
and two learning rates.}
\label{tab:partd_l2_full}
\scriptsize
\begin{ruledtabular}
\begin{tabular}{lccrr}
Architecture & $\lambda$ & $\eta_0$ & Train MSE & Val. MSE \\
$[1,50,1]$ & $0$ & 0.001 & 0.013271 & 0.013980 \\
$[1,50,1]$ & $0$ & 0.01 & 0.009019 & 0.011501 \\
$[1,50,1]$ & $10^{-8}$ & 0.001 & 0.013271 & 0.013980 \\
$[1,50,1]$ & $10^{-8}$ & 0.01 & 0.009019 & 0.011501 \\
$[1,50,1]$ & $10^{-6}$ & 0.001 & 0.013271 & 0.013978 \\
$[1,50,1]$ & $10^{-6}$ & 0.01 & 0.009013 & 0.011509 \\
$[1,50,1]$ & $10^{-4}$ & 0.001 & 0.013265 & 0.013958 \\
$[1,50,1]$ & $10^{-4}$ & 0.01 & 0.008787 & 0.011104 \\
$[1,50,1]$ & $10^{-2}$ & 0.001 & 0.017186 & 0.018577 \\
$[1,50,1]$ & $10^{-2}$ & 0.01 & 0.029600 & 0.025067 \\
$[1,50,1]$ & $10^{-1}$ & 0.001 & 0.027365 & 0.025614 \\
$[1,50,1]$ & $10^{-1}$ & 0.01 & 0.098050 & 0.086916 \\
$[1,100,1]$ & $0$ & 0.001 & 0.012216 & 0.013462 \\
$[1,100,1]$ & $0$ & 0.01 & 0.008833 & 0.011267 \\
$[1,100,1]$ & $10^{-8}$ & 0.001 & 0.012216 & 0.013462 \\
$[1,100,1]$ & $10^{-8}$ & 0.01 & 0.008838 & 0.011282 \\
$[1,100,1]$ & $10^{-6}$ & 0.001 & 0.012216 & 0.013461 \\
$[1,100,1]$ & $10^{-6}$ & 0.01 & 0.008833 & 0.011278 \\
$[1,100,1]$ & $10^{-4}$ & 0.001 & 0.012182 & 0.013358 \\
$[1,100,1]$ & $10^{-4}$ & 0.01 & 0.008459 & 0.011782 \\
$[1,100,1]$ & $10^{-2}$ & 0.001 & 0.016339 & 0.015926 \\
$[1,100,1]$ & $10^{-2}$ & 0.01 & 0.029810 & 0.025455 \\
$[1,100,1]$ & $10^{-1}$ & 0.001 & 0.025476 & 0.023543 \\
$[1,100,1]$ & $10^{-1}$ & 0.01 & 0.098717 & 0.087730 \\
$[1,50,50,1]$ & $0$ & 0.001 & 0.008226 & 0.012321 \\
$[1,50,50,1]$ & $0$ & 0.01 & 0.007618 & 0.011346 \\
$[1,50,50,1]$ & $10^{-8}$ & 0.001 & 0.008256 & 0.012143 \\
$[1,50,50,1]$ & $10^{-8}$ & 0.01 & 0.007651 & 0.011061 \\
$[1,50,50,1]$ & $10^{-6}$ & 0.001 & 0.008496 & 0.011373 \\
$[1,50,50,1]$ & $10^{-6}$ & 0.01 & 0.007880 & 0.011846 \\
$[1,50,50,1]$ & $10^{-4}$ & 0.001 & 0.008531 & 0.011149 \\
$[1,50,50,1]$ & $10^{-4}$ & 0.01 & 0.007904 & 0.011323 \\
$[1,50,50,1]$ & $10^{-2}$ & 0.001 & 0.009965 & 0.010407 \\
$[1,50,50,1]$ & $10^{-2}$ & 0.01 & 0.018444 & 0.016989 \\
$[1,50,50,1]$ & $10^{-1}$ & 0.001 & 0.035086 & 0.031398 \\
$[1,50,50,1]$ & $10^{-1}$ & 0.01 & 0.099508 & 0.088353 \\
$[1,100,100,1]$ & $0$ & 0.001 & 0.007765 & 0.011603 \\
$[1,100,100,1]$ & $0$ & 0.01 & 0.007927 & 0.010980 \\
$[1,100,100,1]$ & $10^{-8}$ & 0.001 & 0.007869 & 0.011428 \\
$[1,100,100,1]$ & $10^{-8}$ & 0.01 & 0.008052 & 0.010607 \\
$[1,100,100,1]$ & $10^{-6}$ & 0.001 & 0.007892 & 0.011752 \\
$[1,100,100,1]$ & $10^{-6}$ & 0.01 & 0.007914 & 0.010750 \\
$[1,100,100,1]$ & $10^{-4}$ & 0.001 & 0.008254 & 0.011307 \\
$[1,100,100,1]$ & $10^{-4}$ & 0.01 & 0.007872 & 0.011625 \\
$[1,100,100,1]$ & $10^{-2}$ & 0.001 & 0.009133 & 0.010249 \\
$[1,100,100,1]$ & $10^{-2}$ & 0.01 & 0.018165 & 0.016952 \\
$[1,100,100,1]$ & $10^{-1}$ & 0.001 & 0.025226 & 0.022559 \\
$[1,100,100,1]$ & $10^{-1}$ & 0.01 & 0.099508 & 0.088352 \\
\end{tabular}
\end{ruledtabular}
\end{table*}


\section{Use of AI/LLM Tools}

\subsection{Tools Used}

ChatGPT (OpenAI, web interface, accessed during the development of
Project~2 through October and November 2026) was used as a conceptual tutor,
programming assistant and drafting/editing tool. No other LLM was used
for the work described here. The author selected the experiments,
executed the programs, inspected the figures and is responsible for the
submitted implementation and scientific interpretation.

\subsection{Writing and Editing}

ChatGPT helped connect explanations from Chapters~3, 5 and 8 of the
course notes \cite{Book_Ch3, Book_Ch5, Book_Ch8} to the implemented algorithms, interpret numerical outputs
and figures, and draft or revise portions of the report. Its substantive
text contributions, including those concerning the MNIST experiments and
the critical evaluation, are identified in Table~\ref{tab:llm_writing}.
Levels denote the extent of assistance; Level~3 indicates substantial
generated draft text subsequently reviewed and edited by the author.

\begin{table*}[t]
\centering
\caption{LLM assistance with report writing and editing.}
\label{tab:llm_writing}
\begin{ruledtabular}
\begin{tabular}{p{0.23\textwidth}cp{0.64\textwidth}}
Section & Level & Nature of assistance \\
Abstract and Introduction & 3 & Drafted a synthesis of the assignment,
methods and measured results; the author reviewed the wording and factual claims. \\
Theory & 1 & Helped revise transitions and notation against the supplied lecture notes. \\
Methods and Implementation & 3 & Drafted methodological descriptions from
the author's code and experimental settings, including dataset splits,
optimization, classification, regularization and early stopping. \\
Results and Discussion & 2 & Helped interpret
learning curves, hyperparameter comparisons and confusion matrix, and drafted
the results tables without inventing measured values. \\
Critical Evaluation and Conclusions & 1 & Wrote the comparative synthesis
of the regression and classification algorithms, strengths, limitations
and follow-up experiments from the authors draft. \\
AI/LLM Declaration & 2 & Reorganized and redrafted this declaration from
the conversation history and the author's code/report, retaining the
course's requested structure. \\
\end{tabular}
\end{ruledtabular}
\end{table*}

The author supplied guidelines and theory basis, relevant lecture chapters,
Python source files, validation/test outputs and generated PDFs as the
basis for these exchanges. Suggested prose was used only after review
for correctness, consistency with the report and agreement with the
actual program output. The LLM did not run the author's experiments
or supply fabricated performance measurements.

\subsection{Programming Assistance}

ChatGPT provided some example code, especially
for adapting the existing regression implementation to classification.
Its suggestions included debugging recording of training histories,
adding loss and accuracy tracking, checking analytical gradients with
finite differences, testing architecture/activation/penalty grids,
implementing checkpoint-based early stopping and creating figures.
The principal file-level contributions are summarized in
Table~\ref{tab:llm_programming}.

\begin{table*}[t]
\centering
\caption{LLM assistance with the submitted Python source.}
\label{tab:llm_programming}
\begin{ruledtabular}
\begin{tabular}{p{0.18\textwidth}cp{0.69\textwidth}}
File & Level & Nature of assistance \\
\texttt{main.py} & 1 & Supplied example of hyperparameter grids for
Runge and MNIST, scikit-learn references and confusion-matrix analysis;
the author incorporated, adapted and executed the suggestions. \\
\texttt{utilities.py} & 2 & Assisted in implementing and extending
optimization and backpropagation, L1/L2 gradients, stable softmax,
one-hot encoding, cross-entropy, accuracy, finite-difference checks
and early-stopping checkpoints. \\
\texttt{plot\_generator.py} & 3 & Supplied or revised plotting routines
for regression learning curves, architecture comparisons, MNIST
learning curves/heatmaps and the confusion matrix; the figures were
produced and reviewed by the author. \\
\end{tabular}
\end{ruledtabular}
\end{table*}

% ============================================================

\bibliographystyle{apsrev4-1}
\bibliography{References}
%TC:endignore
\end{document}
